% Part: computability
% Chapter: computability-theory
% Section: introduction

\documentclass[../../../include/open-logic-section]{subfiles}

\begin{document}

% SEGMENT OLP-0229-S01

\olfileid{cmp}{thy}{int}
\olsection{அறிமுகம்}

\emph{கணிப்புத்தன்மைக் கோட்பாடு} என்று அறியப்படும் தருக்கவியலின்
கிளை, சார்புகள் மற்றும் கணங்களின் கணிப்புத்தன்மை அல்லது சார்புக்
கணிப்புத்தன்மை குறித்த சிக்கல்களைக் கையாள்கிறது. இத்துறை முன்பு
\emph{மீள்வரையறைக் கோட்பாடு} என்று அறியப்பட்டதும், இன்று இரு
பெயர்களுமே பொதுவாகப் பயன்படுத்தப்படுவதும் க்ளீனியின் செல்வாக்குக்குச்
சான்று.

ஏதேனும் கணக்கீட்டு மாதிரியில் கணிக்கக்கூடிய சார்பு
~$f\colon \Nat \pto \Nat$ ஐ \emph{பகுதியாகக் கணிக்கத்தக்கது}
என அழைப்போம். $f$ முழுச் சார்பு எனில், $f$ ஆனது
\emph{கணிக்கத்தக்கது} என்று மட்டும் கூறுவோம். கணிக்கத்தக்க
சிறப்பியல்புச் சார்பு~$\Char{R}$ உடைய உறவு~$R$ உம்
கணிக்கத்தக்கது எனப்படும். $f$ மற்றும்~$g$ பகுதிச் சார்புகள் எனில்,
$f$ ஆனது~$x$ இல் வரையறுக்கப்பட்டுள்ளது, அதாவது~$x$ ஆனது~$f$ இன்
சார்பகத்தில் உள்ளது, என்பதைக் குறிக்க $f(x) \fdefined$ என எழுதுவோம்;
எதிர்மாறாக, $f$ ஆனது~$x$ இல் வரையறுக்கப்படவில்லை என்பதைக் குறிக்க
$f(x) \fundefined$ என எழுதுவோம். $f(x)$ மற்றும் $g(x)$ இரண்டும்
வரையறுக்கப்படாமல் இருக்கின்றன, அல்லது இரண்டும் வரையறுக்கப்பட்டுச்
சமமாக உள்ளன, என்பதைக் குறிக்க $f(x) \simeq g(x)$ ஐப் பயன்படுத்துவோம்.

குறிப்பிட்ட ஒரு கணக்கீட்டு மாதிரியைச் சுட்டாமல் கணிப்புத்தன்மைக்
கோட்பாட்டை ஆராயலாம். இதைச் செய்ய, அனைத்தளாவிய பகுதியாகக் கணிக்கத்தக்க
சார்பு~$\fn{Un}(k, x)$ ஒன்று உள்ளது என்று காட்ட வேண்டும்.
இது பகுதியாகக் கணிக்கத்தக்க சார்புகளை எண்ணிட அனுமதிக்கிறது.
$k$-ஆவது ஓரிட பகுதியாகக் கணிக்கத்தக்க சார்பைக் குறிக்க
~$\cfind{k}$ என்ற குறிமுறையை ஏற்போம்; அது
$\cfind{k}(x) \simeq \fn{Un}(k, x)$ ஆல் வரையறுக்கப்படுகிறது.
(இந்நோக்கத்திற்காக க்ளீனி $\{ k \}$ ஐப் பயன்படுத்தினார்; ஆனால்
அண்மையில் இக்குறிமுறை அவ்வளவாகப் பயன்படுத்தப்படவில்லை.)
சற்றுப் பொதுவாக, தன்னிச்சையான இடஎண்களுடைய பகுதியாகக் கணிக்கத்தக்க
சார்புகளை ஒரே சீராக எண்ணிடலாம்; $k$-ஆவது $n$-இட பகுதி
மீள்வரையறைச் சார்பைக் குறிக்க $\cfind{k}[n]$ ஐப் பயன்படுத்துவோம்.

$f(\vec x, y)$ ஒரு முழுச் சார்பாகவோ பகுதிச் சார்பாகவோ இருந்தால்,
~$f(\vec x, 0)$, \dots, $f(\vec x, y-1)$ அனைத்தும்
வரையறுக்கப்பட்டுள்ளன என்ற நிபந்தனையுடன், $f(\vec x, y) = 0$
ஆகுமாறு உள்ள மிகச்சிறிய~$y$ ஐத் திருப்பித்தரும்~$\vec x$ இன்
சார்பே $\umin{y}{f (\vec x, y)}$; அத்தகைய $y$ எதுவும் இல்லையெனில்,
$\umin{y}{f (\vec x, y)}$ வரையறுக்கப்படவில்லை. $R(\vec x, y)$
ஓர் உறவு எனில், $R(\vec x, y)$ மெய்யாகுமாறு உள்ள மிகச்சிறிய~$y$
ஆக $\umin{y}{R(\vec x, y)}$ வரையறுக்கப்படுகிறது; வேறுவிதமாகக்
கூறினால், $1 \tsub \Char{R}(\vec x, y) = 0$ ஆகுமாறு உள்ள
மிகச்சிறிய~$y$.

ஒரு சார்பு கணிக்கத்தக்கது என்று காட்ட இரு வழிகளில் செயல்படலாம்:
\begin{enumerate}
\item துல்லியமாக: ஒரு டியூரிங் பொறி அல்லது பகுதி மீள்வரையறைச்
  சார்பை வெளிப்படையாக விவரித்து, மனத்தில் கொண்டுள்ள சார்பை அது
  கணிக்கிறது என்று காட்டுக;
\item முறைசாராமல்: அதைக் கணிக்கும் வழிமுறையை விவரித்து,
  சர்ச்சின் முன்வைப்பைப் பயன்படுத்துக.
\end{enumerate}
இவ்விரண்டிற்கும் இடையில் கூர்மையான எல்லை இல்லை; ஒரு வழிமுறையின்
விரிவான விளக்கம், கோட்பாட்டளவில் உரிய டியூரிங் பொறியையோ பகுதி
மீள்வரையறை வரையறைகளின் தொடரையோ எவ்வாறு வடிவமைக்கலாம் என்பது
ஓரளவு தெளிவாக இருக்கப் போதுமான தகவலைத் தர வேண்டும். முழுத்
துல்லியமான வரையறைகள் தகவலளிப்பவையாக இருக்க வாய்ப்பில்லை; எனவே
இவ்விரு அணுகுமுறைகளுக்கும் இடையில் பொருத்தமான நடுவழியைக் காண
முயல்வோம். அதாவது, துல்லியமான வரையறைகளைப் பெற முடியும் என்பதற்கான
உள்ளுணர்வான, அதேவேளை துல்லியமான நிரூபிப்புகளைக் காண முயல்வோம்.

\end{document}
